\chapter{Background}
\label{chap:technical}

This chapter introduces the technical background necessary to understand
the reproduction experiments and architectural extension described in later
chapters. The content is organised to follow a clear conceptual progression:
starting from the definition of semantic segmentation and the architectural
components used to solve it, then moving to the problem of learning under
limited annotation, the role of self-supervised depth estimation in
addressing that problem, the multi-task architecture that combines
segmentation and depth, and finally the motivation for adding an explicit
cross-task consistency objective. Each section prepares the ground for the next and culminates in a
summary of the research gap addressed in Chapter~\ref{chap:execution}.

% -----------------------------------------------------------------------------
\section{Semantic Segmentation}
\label{sec:semseg}

Semantic segmentation is a dense prediction task that assigns a semantic
class label to every pixel in an image. Unlike image-level classification,
which produces a single category label for the entire input, semantic
segmentation must produce a categorical decision at every spatial location,
yielding a prediction map with the same height and width as the input
image. This pixel-level understanding is essential in applications such as
autonomous driving, where a vehicle must distinguish drivable road from
pedestrians, obstacles, and static infrastructure at every point in the
scene~\cite{cordts2016cityscapes}.

Segmentation performance is measured by mean Intersection-over-Union
(mIoU). For each class~$c$, IoU computes the ratio of overlap between the
predicted region and the ground-truth region:
\[
    \mathrm{IoU}_c = \frac{TP_c}{TP_c + FP_c + FN_c},
\]
where $TP_c$, $FP_c$, and $FN_c$ are the numbers of true positive, false
positive, and false negative pixels for class~$c$ respectively. The final
mIoU is the arithmetic mean across all classes. mIoU is preferred over
simple pixel accuracy because class frequencies in real scenes are highly
imbalanced: road and sky pixels far outnumber pedestrians or traffic
lights. By computing a separate score for each class before averaging,
mIoU ensures that rare but semantically important classes contribute
equally to the overall measure.

The high cost of annotation is a defining practical constraint of this
task. Labelling a single complex urban image at the pixel level requires
a human annotator to carefully outline every visible object, a process
that takes around 1.5~hours per image on average~\cite{cordts2016cityscapes}.
This cost makes it infeasible to annotate large datasets in full, and
motivates research into training strong segmentation models from a small
number of labelled examples. The present project addresses this
label-efficient problem directly, and the experimental setup---including
the Cityscapes benchmark and the specific annotation budgets
used---is described in Chapter~\ref{chap:execution}.

% -----------------------------------------------------------------------------
\section{Deep Architectures for Dense Prediction}
\label{sec:architectures}

Effective semantic segmentation requires architectural components that can
both extract rich semantic representations from images and recover the
spatial precision lost during feature extraction. This section describes the building blocks that collectively form the
segmentation and depth architectures reproduced and extended in this
project.

\subsection{Encoder--Decoder Networks}

Modern segmentation models follow an encoder--decoder
structure~\cite{long2015fcn}. The encoder, typically a deep convolutional
backbone, progressively reduces spatial resolution through pooling and
strided convolutions. As depth increases, each feature map captures
increasingly abstract semantic information over a larger receptive field,
but at the cost of spatial detail. The decoder reverses this process,
gradually restoring spatial resolution so that the final output matches
the input size. This structure is fundamental to dense prediction: the
encoder builds semantic understanding, while the decoder projects that
understanding back onto the pixel grid.

\subsection{ResNet-101 Backbone}

The reproduced framework uses ResNet-101~\cite{he2016resnet} as its
encoder backbone.
Standard deep networks suffer from vanishing gradients during
backpropagation: as gradient signals travel through many layers, repeated
multiplication by small derivatives causes them to diminish exponentially,
leaving early layers with almost no useful update signal. ResNets resolve
this by adding a shortcut connection that bypasses each residual block:
\[
    \mathbf{y} = \mathcal{F}(\mathbf{x}) + \mathbf{x},
\]
where $\mathbf{x}$ is the block input, $\mathcal{F}(\mathbf{x})$
represents the convolutional path, and $\mathbf{y}$ is the block output.
The shortcut provides a direct gradient path from any layer back to
the input, enabling stable training of very deep networks. With 101
layers, ResNet-101 provides strong general-purpose visual features and
is a standard encoder choice in segmentation frameworks. The encoder is
initialised with weights pre-trained on
ImageNet~\cite{deng2009imagenet}, providing a strong starting point that
generalises well to downstream tasks, especially when labelled data is limited.

\subsection{Multi-Scale Context with ASPP}

A recurring challenge in segmentation is that objects vary greatly in
apparent size depending on their distance from the camera. A module that
only aggregates local context cannot simultaneously handle distant
pedestrians (appearing small) and nearby buildings (appearing large).
Atrous Spatial Pyramid Pooling
(ASPP)~\cite{chen2018deeplabv3plus} addresses this by running multiple
atrous (dilated) convolutions in parallel at different dilation rates.
Atrous convolutions insert gaps between kernel elements, expanding the
effective receptive field without increasing parameter count or reducing
spatial resolution. By pooling responses at several scales and
concatenating them, ASPP captures multi-scale context from a single
feature map. The reproduced framework uses ASPP at the encoder bottleneck before
decoder upsampling; implementation details are given in
Chapter~\ref{chap:execution}.

\subsection{U-Net Decoder and Skip Connections}

Upsampling from the low-resolution bottleneck to a full-resolution
prediction map with bilinear interpolation alone is insufficient for
accurate boundary prediction: the fine spatial detail has been discarded
by the encoder's downsampling steps and cannot be recovered from the
deep features alone. U-Net~\cite{ronneberger2015unet} addresses this
with skip connections that concatenate the encoder's feature maps at each
resolution directly onto the corresponding decoder layer. This gives the
decoder simultaneous access to high-level semantic context from deep
layers and low-level edge detail from shallow layers, producing
predictions that are both semantically accurate and spatially precise.

The reproduced framework uses this ASPP-plus-decoder design for both the
segmentation and depth branches, with both branches sharing an identical
decoder architecture. This design consistency enables decoder weights
trained on the depth task to initialise the segmentation decoder
directly, transferring geometric priors without architectural adaptation.
This role of SDE pre-training is discussed further in
Section~\ref{sec:sde}. Implementation details are given in
Chapter~\ref{chap:execution}.

\subsection{Batch Normalisation}

Training deep networks is stabilised by Batch Normalisation
(BN)~\cite{ioffe2015batchnorm}, which normalises layer activations across
the training batch at each step. This reduces sensitivity to weight
initialisation and allows higher learning rates. In multi-GPU setups,
Synchronised BN (SyncBN) aggregates statistics across all devices so
that normalisation is performed over the full effective batch. This
distinction is directly relevant to the reproduced experiments: all
runs were conducted on a single GPU with batch size~2, so SyncBN could
not provide cross-device statistic aggregation.
The practical consequences of this constraint are discussed in
Chapter~\ref{chap:execution}.

% -----------------------------------------------------------------------------
\section{Label-Efficient Semantic Segmentation}
\label{sec:ssl}

\subsection{The Label-Efficient Problem}

The label-efficient setting studies how to train strong models when only
a small fraction of training images carry pixel-level annotations. The
remaining images are available but unlabelled. This is distinct from the
fully supervised setting, where every training image is annotated, and
from the unsupervised setting, where no labels exist at all. The goal is
to close the gap between limited-supervision and full-supervision
performance by exploiting the structural information in unlabelled data.

Semi-supervised learning (SSL) is the dominant approach to this
problem~\cite{chapelle2006ssl}. In the segmentation context, the model
is trained on a small labelled subset of images alongside a much larger
unlabelled pool. The supervised objective is a pixel-wise cross-entropy
loss over labelled images; the semi-supervised objective encourages the
model to extract useful representations from the unlabelled ones. The
formal problem setup, including notation and the specific annotation
budgets used in this project, is given in Chapter~\ref{chap:execution}.

\subsection{Pseudo-Labels and Confirmation Bias}

Pseudo-labelling is the most widely used SSL strategy for
segmentation~\cite{lee2013pseudolabel}. The model generates predictions
for unlabelled images, treats those predictions as ground-truth labels,
and trains on the combined labelled and pseudo-labelled data. The process
can be iterated as predictions improve.

However, pseudo-labelling introduces a fundamental risk: confirmation
bias~\cite{arazo2020pseudolabel}. When the model's initial performance is
poor, the pseudo-labels it generates contain many errors. Training on
those incorrect labels reinforces the model's mistakes, causing
performance to stagnate or degrade. This risk is most severe when very
few labelled images are available, since the initial pseudo-labels are
least reliable precisely when the most leverage from unlabelled data is
needed.

\subsection{Mean Teacher}

The Mean Teacher framework~\cite{tarvainen2017meanteacher} helps reduce
confirmation bias by separating the pseudo-label generator from
the model being trained. Two networks are maintained: a student, updated
by standard gradient descent, and a teacher, whose weights are an
exponential moving average (EMA) of the student's weights:
\[
    \theta_T \leftarrow \alpha\,\theta_T + (1 - \alpha)\,\theta_S,
\]
where $\theta_T$ and $\theta_S$ are the teacher and student weights
respectively, and $\alpha$ is a smoothing coefficient close to~1.
Because the teacher accumulates a running
average of the student's entire training history, it behaves as a
temporal ensemble and produces more stable, consistent predictions than
any single student checkpoint. These stable predictions serve as pseudo-labels for unlabelled data,
helping to reduce the noisy feedback that causes confirmation bias.

\subsection{Consistency Regularisation and Data Augmentation}

Consistency regularisation complements pseudo-labelling by requiring the
model to produce the same prediction for a given image regardless of what
perturbations are applied to it~\cite{sohn2020fixmatch}. This encourages
the model to learn representations that are genuinely semantic rather
than tied to surface-level visual cues.

Within the Mean Teacher framework, a \emph{weak--strong augmentation
pairing} strategy implements this principle. For each unlabelled image,
two versions are created: a weakly augmented version with mild
transformations is passed to the teacher to generate reliable pseudo-labels;
a strongly augmented version (additionally applying colour jitter, random
greyscale conversion, and Gaussian blur) is passed to the student. The
student is trained to match the teacher's pseudo-labels despite the more
challenging input, encouraging it to learn features that are robust to
appearance variation.

\subsection{Data Mixing Strategies}
\label{sec:data-mixing}

Data mixing is a complementary augmentation strategy that creates
synthetic training samples by combining pixels from two images according
to a binary mask. ClassMix~\cite{olsson2021classmix} builds the mixing
mask from class segments, forcing the model to recognise the same class
in varied contexts. However, ClassMix ignores scene geometry and can
produce physically implausible composites---for example, pasting a
distant building in front of a nearby car.

DepthMix~\cite{hoyer2023ijcv_sde} resolves this by using the predicted
depth ordering between two images to decide which regions should appear
in front in the mixed sample. In effect, closer regions are allowed to
occlude more distant regions, consistent with real-world geometry. This
produces synthetic samples that respect depth ordering and better support
learning accurate object boundaries. The depth estimates required for
DepthMix come from the self-supervised depth estimation component
described in the following section.

% -----------------------------------------------------------------------------
\section{Self-Supervised Depth Estimation for Segmentation}
\label{sec:sde}

\subsection{Motivation: Depth as a Geometric Prior}

Depth and semantics are closely related in natural scenes. Objects within
the same semantic category tend to occupy consistent depth ranges, and
sharp depth discontinuities frequently coincide with semantic boundaries:
the edge of a pedestrian against a building background is simultaneously
a semantic boundary and a depth boundary. Incorporating geometric
information from depth can therefore make segmentation features more
spatially precise and structurally coherent.

Obtaining ground-truth depth annotations, however, requires expensive
equipment such as LiDAR sensors and does not scale to large datasets.
Self-supervised depth estimation (SDE) offers an alternative: it learns
to predict depth from unlabelled video sequences by exploiting the
geometric consistency between consecutive frames, requiring no manual
annotation~\cite{godard2017unsupervised}.

\subsection{View Synthesis and Photometric Loss}

SDE trains a depth network and a pose network jointly. Given consecutive
frames from a video, the depth network predicts a dense depth map for
each frame, and the pose network estimates the camera motion between
frames. The source frame is then warped into the viewpoint of the target
frame using the predicted depth and pose. If the depth and pose
predictions are accurate, the synthesised view should closely match the
real target frame. Training minimises the photometric reprojection error
between synthesised and real target frames, combining an $\ell_1$
pixel-intensity term with a Structural Similarity
(SSIM)~\cite{wang2004ssim} term. Because the entire supervision signal comes from the images themselves,
SDE can be trained on large volumes of unlabelled driving video, including
video sequences that commonly accompany autonomous-driving datasets,
building up rich geometric priors at no annotation cost.

\subsection{Depth Pre-training Variants: \texttt{fd0} and \texttt{fd2}}
\label{sec:fd0-fd2}

Within the framework of Hoyer et al.~\cite{hoyer2023ijcv_sde}, SDE
pre-training produces two encoder weight variants that differ in how the
encoder is regularised during depth pre-training.

\textbf{\texttt{fd0}} applies standard SDE pre-training, optimising the
encoder purely to minimise the photometric loss. While this produces an
encoder with strong geometric representations, it risks drifting away
from the ImageNet feature space: the SDE objective provides no incentive
to retain general visual features, so the encoder may lose the semantic
knowledge acquired during ImageNet pre-training. When these weights
initialise the segmentation branch, the lost semantic features can be
difficult to recover.

\textbf{\texttt{fd2}} addresses this by adding a feature-distance
regularisation term~$\mathcal{L}_F$ during depth
pre-training~\cite{hoyer2023ijcv_sde}, which penalises the $\ell_2$
distance between the current encoder's features~$f^E_\theta$ and those
of the frozen ImageNet-initialised encoder~$f^E_I$:
\[
    \mathcal{L}_F = \| f^E_\theta - f^E_I \|_2,
    \qquad
    \mathcal{L}_{D,\text{pretrain}} = \mathcal{L}_D + \lambda_F \mathcal{L}_F.
\]
By keeping the encoder close to its ImageNet initialisation, \texttt{fd2}
retains general semantic representations while still acquiring geometric
knowledge, making it a stronger starting point for the downstream
segmentation task. These two variants are used in the reproduction of the
SDE feature transfer experiments discussed in Chapter~\ref{chap:evaluation}.

\subsection{How SDE Benefits Label-Efficient Segmentation}

Hoyer et al.~\cite{hoyer2023ijcv_sde} demonstrate several ways in which
SDE improves label-efficient segmentation, each targeting a
different stage of the training pipeline.

First, \textbf{SDE-guided data selection} uses features from a
pre-trained depth model as a proxy for sample informativeness. Diversity
sampling (DS) selects images that cover a wide range of scene content,
while uncertainty sampling (US) selects images where the depth model is
least confident. Because depth estimation and semantic segmentation share
similar notions of scene complexity, SDE-based selection signals transfer
to the segmentation task without requiring any segmentation labels for
the selection step. In the original framework, combining both criteria
(DS+US) is reported to provide stronger selection performance than using
either criterion alone.

Second, \textbf{DepthMix} (introduced in Section~\ref{sec:data-mixing}) uses
SDE-predicted depth maps to enforce geometric consistency when composing
synthetic training images, producing more physically plausible augmented
samples than class-based mixing.

Third, SDE pre-trained encoder weights can be used to \textbf{initialise
the segmentation encoder} through direct weight transfer, or maintained
as a \textbf{parallel depth task} throughout segmentation training via
multi-task learning. The multi-task approach is described in the
following section. Together, these mechanisms address complementary
bottlenecks in the label-efficient pipeline, and their individual and
combined contributions are quantified in Chapter~\ref{chap:evaluation}.

% -----------------------------------------------------------------------------
\section{Multi-Task Learning and PAD-Based Feature Exchange}
\label{sec:mtl}

\subsection{Multi-Task Learning}

Multi-task learning (MTL) trains a single model on multiple related tasks
simultaneously by sharing representations across
them~\cite{caruana1997multitask}. When semantic segmentation and depth
estimation are trained jointly, the shared encoder receives gradient
signals from both task losses and is guided towards features that are
jointly useful for geometry and semantics. Compared to transfer
learning---where SDE pre-trained weights are used only as an
initialisation---MTL maintains the depth supervision signal throughout
the entire training process. This provides a continuous source of
geometric regularisation that is especially valuable when labelled
segmentation data is scarce, since the depth task can be trained on all
available unlabelled frames without requiring any segmentation annotations.

\subsection{PAD-Net}

The multi-task architecture reproduced and extended in this project is
based on PAD-Net~\cite{xu2018padnet}. PAD-Net maintains two parallel decoder
branches on top of a shared encoder: one for depth estimation and one
for semantic segmentation. Both branches use the ASPP and U-Net
upsampling block structure described in Section~\ref{sec:architectures}.
At an intermediate decoder stage, the two branches exchange information
through an attention-guided distillation module. Each branch applies
self-attention to its own intermediate features, then incorporates the
other branch's attended features via element-wise addition. Schematically:
\[
    \tilde{f}_\text{seg}   = f_\text{seg}   +
        \mathrm{SA}_\text{depth}(f_\text{depth}),
    \qquad
    \tilde{f}_\text{depth} = f_\text{depth} +
        \mathrm{SA}_\text{seg}(f_\text{seg}),
\]
where $\mathrm{SA}(\cdot)$ denotes the self-attention operation,
$f_\text{seg}$ and $f_\text{depth}$ are the intermediate features of
each branch at the distillation layer, and $\tilde{f}$ denotes the
updated features after cross-task exchange. In concrete terms: the
segmentation branch looks at what the depth branch has learned about
scene geometry and absorbs it into its own features, and vice versa. The
updated features $\tilde{f}$ are then passed through the remaining
decoder layers to produce final predictions.

\subsection{Limitations of Implicit Feature Exchange}

While PAD-Net's attention-guided distillation allows each branch to
incorporate features from the other, this interaction is fundamentally
implicit. It modulates attention outputs rather than placing any
direct constraint on the feature representations themselves. There is no
objective that explicitly requires the segmentation and depth decoder
features to be consistent with one another.

A second, more structural issue compounds this. The depth decoder
receives self-supervised training signals from all available unlabelled
frames throughout training, while the segmentation decoder is supervised
only on the small labelled subset. This asymmetry means that the depth
branch may accumulate richer and more geometrically grounded representations
than the segmentation branch, yet PAD-Net provides no direct feature-level
objective to enforce this transfer into the segmentation feature space.
The attention exchange operates symmetrically and does not compensate for
the supervision asymmetry.

These limitations motivate the addition of an explicit cross-task
feature alignment objective, introduced in the following section.

% -----------------------------------------------------------------------------
\section{Cross-Task Consistency and Feature Alignment}
\label{sec:cross-task}

\subsection{The Case for Explicit Feature Alignment}

The supervision asymmetry described in Section~\ref{sec:mtl} suggests a
specific opportunity: if the depth branch has learned richer geometric
representations due to its access to more training signal, then
explicitly encouraging the segmentation decoder to align its features
with the depth decoder's could transfer that geometric knowledge directly
into the segmentation feature space. This alignment does not require any
additional annotations---the depth decoder's features are a continuously
available, label-free source of geometric supervision.

An equally valid question is whether the opposite direction is more
appropriate. The segmentation branch receives direct supervision from
ground-truth labels, making its features semantically grounded. Aligning
the depth branch's features towards the segmentation branch's would
instead treat the labelled segmentation task as the semantic anchor.
Both directions represent different assumptions about which task's
representations should guide the other, and both are considered in the
experimental design described in later chapters.

\subsection{Related Work on Cross-Task Consistency}

The idea that predictions or representations from different tasks should
be mutually consistent has been explored from several angles.
Zamir et al.~\cite{zamir2020crosstask} formalise it through inference-path
invariance: predictions obtained via different combinations of tasks
should remain geometrically consistent regardless of the path taken.
Nakano et al.~\cite{nakano2021crosstask} apply this principle directly to
multi-task learning by constraining intermediate features across task
branches, demonstrating improved performance. Li et
al.~\cite{li2022mtl} study multi-task dense prediction under partial
annotations, showing that cross-task relationships can be exploited when
different tasks have access to different amounts of labelled
data---a situation directly analogous to the label-efficient multi-task
setup considered here.

These works collectively suggest that explicit feature-level constraints
can supplement the implicit sharing that a shared encoder and attention
exchange provide, and that such constraints are particularly valuable
when task supervision is asymmetric.

\subsection{Feature Alignment: Loss Formulations}
\label{sec:ct-losses}

Implementing cross-task consistency requires a loss function that
quantifies the difference between intermediate features from the two
decoder branches without being confounded by differences in their
absolute activation magnitudes, since the depth and segmentation branches
may operate at different numerical scales. Scale-insensitive formulations
that compare the \emph{direction} of feature vectors rather than their
absolute values are therefore appropriate.

Two such objectives are considered in this project: normalised mean
squared error (normalised MSE), which compares $\ell_2$-normalised
feature vectors, and cosine similarity loss, which directly penalises
angular disagreement between feature directions. Both avoid comparing
raw activation magnitudes and are naturally suited to aligning
representations across task-specific decoder branches. Their exact
formulations and integration into the proposed $\mathcal{L}_\text{ct}$
are given in Chapter~\ref{chap:execution}.

% -----------------------------------------------------------------------------
\section{Summary and Research Gap}
\label{sec:summary-gap}

This chapter has introduced the technical foundations of this project
along a connected conceptual arc. Semantic segmentation is a dense
prediction task where annotation cost makes fully supervised training
impractical at scale. Semi-supervised approaches address this by
exploiting unlabelled data through pseudo-labelling and consistency
regularisation, with the Mean Teacher framework providing stable pseudo-labels that help
mitigate confirmation bias. Self-supervised depth estimation offers a
complementary route: by training on freely available
unlabelled video, it produces geometric priors that can improve
label-efficient segmentation through data selection, geometry-aware
mixing, weight transfer, and multi-task learning. The PAD-Net
architecture realises this multi-task combination, allowing the depth
and segmentation decoders to exchange information through attention-guided
distillation.

Two related limitations remain, however. First, PAD-Net's feature
exchange is implicit: the attention mechanism modulates representations
without placing any direct constraint on their consistency. Second, the
depth decoder receives self-supervised signals from all available
unlabelled frames throughout training, while the segmentation decoder is
confined to the small labelled subset. This supervision asymmetry means
the depth branch may accumulate richer geometric representations, but
PAD-Net provides no direct feature-level objective to enforce this
transfer into the segmentation feature space.

This leaves a gap between the implicit cross-task feature exchange
achievable through PAD-Net and the explicit cross-task feature
consistency that a direct alignment objective could enforce. The
following chapter therefore introduces and formalises the proposed
Cross-Task Consistency Loss $\mathcal{L}_\text{ct}$, which adds an
explicit feature-level alignment objective to the PAD-MTL training
pipeline and investigates whether it can close this gap under
label-efficient conditions.